\section{관련 연구}
\label{sec:related}

\subsection{배틀로얄과 e스포츠의 결과 예측}
\label{sec:related:esports}

% 경기 종료 후의 순위 추정과 경기 중 예측의 구분
PUBG의 결과 예측 연구는 경기 종료 후의 요약 통계로 최종 순위를 추정하는 접근에서 출발하였으며~\cite{KagglePUBG}, 같은 유형의 통계로 순위를 예측할 때 유효한 요인과 모델이 경기 모드에 따라 달라짐이 보고되었다~\cite{LeeLee25}.
이러한 접근은 입력이 경기 종료 시점에 완성되므로, 경기 도중에 이용 가능한 정보로 이후의 위험을 예측하는 과제와는 구별된다.
다른 장르의 e스포츠에서는 경기 중에 갱신되는 정보를 이용한 예측이 연구되었다.
MOBA 장르에서는 실시간 승리 예측~\cite{Hodge21}과 시간 구간별 팀 간 상대적 우위의 평가~\cite{Yu19}가 제안되었고, 개인 단위의 짧은 구간 사망 예측도 Dota~2~\cite{Katona19,Ringer23}와 Counter-Strike~\cite{Marshall22}에서 다루어졌다.

% 배틀로얄의 가까운 선행연구와 본 연구의 비교 범위
배틀로얄의 경기 중 예측으로는 Sagong 등이 단계별 활동 및 교전 점수와 플레이 유형으로 개별 플레이어의 생존을 예측하였고~\cite{Sagong25a}, 단계별 전략 시퀀스로 승리를 예측하였다~\cite{Sagong25b}.
KRAFTON의 예측기는 팀의 행동, 상태, 공간 정보를 팀 간 Transformer에 입력하고, 생존분석을 이용하여 팀 단위의 실시간 우승확률을 추정한다~\cite{KRAFTON}.
공개 구현에는 이산시간 hazard(LogisticHazard)를 포함한 여러 생존모형 기준선과 팀 간 attention을 제거한 변형이 포함되어 있다~\cite{KRAFTON}.
이들 연구는 최종 순위와 우승확률, 개인의 생존 여부 등 서로 다른 분석 단위와 예측 구간을 다룬다.
본 연구는 안전구역의 단계 경계를 예측 시점으로 하고, 현재 생존한 팀이 이어지는 한 단계 안에 완전히 탈락할 조건부 확률에 초점을 맞춘다.
이를 위해 활동 점수 기반 설계와 상대위험 기반 설계를 동일한 팀--단계 표본과 평가 조건에 맞추어 적용하고, 입력 표현, 학습 목적, 단계별 기저위험, 모델 구조가 예측 성능에 미치는 영향을 구분하여 평가한다.

\subsection{생존분석과 생태학적 위험 모델링}
\label{sec:related:survival}

% 상대위험, 절대적 사건확률 및 이산시간 모형
생존분석은 사건이 일어나기까지의 시간을 분석하면서, 사건이 관측되기 전에 추적이 종료되는 검열을 함께 다룬다.
Cox 비례위험모형은 기저위험의 함수 형태를 사전에 특정하지 않고 부분우도를 이용하여 공변량의 상대위험을 추정한다~\cite{Cox72}.
상대위험 점수 자체는 사건확률이 아니지만, 추정한 기저위험을 결합하면 생존확률과 구간별 조건부 사건확률을 산출할 수 있다.
이산시간 생존모형은 시간을 구간으로 나누고, 각 구간의 시작까지 생존한 대상이 해당 구간 안에 사건을 겪을 조건부 확률, 즉 구간별 hazard를 모형화한다~\cite{SingerWillett93}.
이때 구간별 기저위험을 통해 시간에 따른 공통적인 위험 수준의 변화를 반영할 수 있다.
신경망 기반 생존모형에서도 구간별 hazard를 학습하는 접근~\cite{Gensheimer19,Ren19}과 사건시간의 확률분포를 학습하는 접근~\cite{Lee18}이 제안되었으며, 시간에 따라 갱신되는 관측 이력을 이용한 동적 예측으로 확장되었다~\cite{Lee20}.

% 단계 경계에서의 동적 예측과 학습 목적 비교
예측 시점마다 그 시점의 생존자를 대상으로 당시까지 관측된 정보를 이용하여 이후의 위험을 추정하는 landmarking은 이러한 동적 예측을 구성하는 방법이다~\cite{vanHouwelingen07}.
본 연구는 각 단계 경계를 landmark로 삼아 관측 이력을 갱신하고, 이어지는 한 단계의 팀 탈락확률을 이산시간 hazard로 추정한다.
따라서 비교의 초점은 Cox 모형이 구간별 위험을 예측할 수 있는지 여부가 아니라, 상대위험 점수를 학습하는 목적과 다음 단계 탈락확률을 직접 학습하는 목적이 동일한 예측 과제에서 어떤 차이를 보이는지에 있다.
이를 위해 입력 표현과 공통 모델 구조를 고정한 상태에서 학습 목적을 비교하고, 단계별 기저위험의 기여를 별도로 평가한다.
Cox 계열 목적으로 학습한 모델에도 학습 집합에서 추정한 단계별 기저위험을 결합하여~\cite{PrenticeGloeckler78}, 기저위험의 유무가 아니라 학습 목적 자체의 차이를 비교한다.

% 생태학적 관점이 공변량 구성에 제공하는 근거와 적용 범위
생태학에서는 서식지의 양과 공간적 배치가 개체군의 존속과 어떻게 연결되는지를 다루어 왔다~\cite{Fahrig03,Hanski94}.
동적 점유 모형은 이전 기간에 점유된 장소에서 다음 기간에 종이 사라질 조건부 확률을 추정하며, 정착과 불완전한 탐지도 함께 모형화한다~\cite{MacKenzie03}.
이러한 접근은 환경 조건에 따른 기간별 존속 상태의 전이를 다룬다는 점에서 팀 탈락 문제와 연결되지만, 고정된 장소의 점유와 재정착을 다루는 모형을 이동하는 팀에 그대로 적용할 수는 없다.
또한 이동생태학은 이동을 개체의 상태와 능력, 외부 환경이 함께 작용하는 과정으로 해석한다~\cite{Nathan08}.
본 연구는 이들 관점을 참고하여 안전구역의 축소 아래에서 팀의 존속과 관련될 수 있는 구역 경계에 대한 노출, 이동 부담, 주변 경쟁 밀도, 남은 팀원 수를 서식지 맥락으로 구성한다.
생태학적 관점은 위험 공변량을 구성하고 그 관계를 해석하는 근거로 사용하며, 공변량과 탈락 위험의 관련성 및 그 일관성은 경기 자료에서 검증한다.

\subsection{시공간 다중 에이전트 모델링}
\label{sec:related:multiagent}

% 상호작용과 시간적 이력을 표현하는 구조
시공간 데이터 마이닝은 이동 궤적과 사건을 이용한 예측, 군집, 이상 탐지 등을 다루어 왔으며, 딥러닝을 이용하여 시공간 의존성을 학습하는 방법도 연구되었다~\cite{Atluri18,Wang22}.
다수의 에이전트가 상호작용하는 장면에서는 에이전트 축과 시간 축의 attention을 분리하여 결합하고, 이를 통해 에이전트 사이의 관계와 시간에 따른 변화를 함께 표현하는 방법이 제안되었다~\cite{Ngiam22,Capellera24}.
그래프 구조를 attention에 반영하는 접근으로는 노드 간 최단경로 거리와 간선 정보를 attention의 편향으로 인코딩하는 Graphormer가 있다~\cite{Ying21}.
스포츠와 e스포츠에서도 선수 간 관계를 그래프로 표현하여 미식축구의 전진 거리와 Counter-Strike의 라운드 승리와 같은 경기 내 결과를 예측한 연구가 있다~\cite{Xenopoulos21}.

% 서식지 표현과 hazard 목적 위에서의 구조적 추가 이득
이러한 구조는 에이전트 간 상호작용과 시간적 이력을 학습할 수 있는 표현력을 제공하지만, 입력 정보, 학습 목적, 모델 구조가 동시에 달라지는 비교에서는 관측된 성능 차이를 구조 자체의 기여로 해석하기 어렵다.
본 연구는 팀 축 attention, 단계 축 attention, 관계 편향을 갖춘 모델을 참조 모델로 사용하되, 서식지 맥락 표현과 단계 조건부 hazard 목적을 고정한 상태에서 각 구조 요소를 한 번에 하나씩 추가하여 그 조건부 기여를 평가한다.
